\subsection{Linear Equations}
\hrulefill

\paragraph{Definition} A \textbf{linear equation} in the variables $x_1, x_2, \ldots, x_n$ is an equation that can be written in the form
\[a_1 x_1 + a_2 x_2 + \cdots + a_n x_n = b\]
where $a_1, a_2, \ldots, a_n$ and $b$ are constants. The constants $a_1, a_2, \ldots, a_n$ are called the \textbf{coefficients} of the equation, and $b$ is called the \textbf{constant term}.

\paragraph*{Examples} of linear equations include:
\begin{itemize}
    \item $2x + 3y - z = 5$
    \item $-x_1 + 4x_2 + 0x_3 = 7$
    \item $0x + 0y + 0z = 0$
\end{itemize}

\paragraph*{Non-examples} of linear equations include:
\begin{itemize}
    \item $x^2 + y = 3$ (because of the $x^2$ term)
    \item $xy + z = 4$ (because of the $xy$ term)
    \item $\sin(x) + y = 1$ (because of the $\sin(x)$ term)
    \item $x + \sqrt{y} = 2$ (because of the $\sqrt{y}$ term)
\end{itemize}

\subsubsection*{Systems of Linear Equations}
\paragraph*{Definition} A \textbf{finite} collection of linear equations involving the same set of variables is called a \textbf{system of linear equations}.
\paragraph*{Example} of a system of linear equations:
\[\begin{cases}
2x + 3y - z = 5 \\
-x + 4y + 2z = 7 \\
3x - y + z = 2 \\
\cdots
\end{cases}\]

\paragraph*{Definition} A sequence of numbers $(s_1, s_2, \ldots, s_n)$ is called a \textbf{solution} to the system of linear equations if, when we substitute $s_i$ for $x_i$ in each equation, the equation holds true.

\paragraph*{Example} Consider the system of equations:
\[\begin{cases}
    x + y + z = 3 \\
    2x + y + 3z = 1 \\
\end{cases}\]
Is the sequence $(-2, -5, 0)$ a solution to this system? Does $x = y = z = 1$?
\paragraph*{Solution} To check if $(-2, -5, 0)$ is a solution, substitute $x = -2$, $y = -5$, and $z = 0$ into both equations:
\begin{itemize}
    \item First equation: $(-2) + (-5) + 0 = -7 \neq 3$ (does not hold)
    \item Second equation: $2(-2) + (-5) + 3(0) = -4 - 5 = -9 \neq 1$ (does not hold)
\end{itemize}
Thus, $(-2, -5, 0)$ is not a solution to the system. \\
Next, check if $(1, 1, 1)$ is a solution:
\begin{itemize}
    \item First equation: $1 + 1 + 1 = 3$ (holds)
    \item Second equation: $2(1) + 1 + 3(1) = 2 + 1 + 3 = 6 \neq 1$ (does not hold)
\end{itemize}
Thus, $(1, 1, 1)$ is also not a solution to the system.

\subsubsection*{Solutions to Linear Systems}
\paragraph*{Definition} There are three possible scenarios for the solutions to a system of linear equations:
\begin{itemize}
    \item \textbf{No solution}: The system is inconsistent, meaning there are no values for the variables that satisfy all equations simultaneously.
    \item \textbf{Unique solution}: There is exactly one set of values for the variables that satisfies all equations.
    \item \textbf{Infinitely many solutions}: There are multiple sets of values for the variables that satisfy all equations. (i.e. the same line or scalar multiples of each other)
\end{itemize}

Systems of linear equations that have at least one solution are called \textbf{consistent}, while those with no solutions are called \textbf{inconsistent}.

\hrulefill

\subsubsection*{Matricies}

\paragraph*{Definition} A \textbf{matrix} is a rectangular array of numbers arranged in rows and columns. The numbers in the matrix are called its \textbf{entries} or \textbf{elements}. 
A matrix with $m$ rows and $n$ columns is called an $m \times n$ matrix.

\paragraph*{Example} of an \textbf{augmented} matrix:
\[\begin{cases}
3x_1 + 2x_2 - x_3 + x_4 = -1 \\
2x_1 - x_3 + 2x_4 = 0 \\
3x_1 + x_2 + 2x_3 + 5x_4 = 2\\
\end{cases} = \begin{bmatrix}
3 & 2 & -1 & 1 & | & -1 \\
2 & 0 & -1 & 2 & | & 0 \\
3 & 1 & 2 & 5 & | & 2 \\
\end{bmatrix}\]

\subsubsection*{Elementary Operations}
\paragraph*{Definition} There are three types of \textbf{elementary row operations} that can be performed on a matrix:
\begin{itemize}
    \item \textbf{Row swapping}: Interchanging two rows of the matrix.
    \item \textbf{Row scaling}: Multiplying all entries in a row by a non-zero constant.
    \item \textbf{Row addition}: Adding a multiple of one row to another row.
\end{itemize}

These operations are used in methods such as Gaussian elimination to solve systems of linear equations.
Row operations do not produce equal matricies, but they do produce \textbf{equivalent} matrices, meaning they represent systems of equations with the same solution set.
In order for 2 matricies to be equal, they must have the same dimensions and identical entries in corresponding positions.

\paragraph*{Example} Consider the matricies:
\[\begin{bmatrix}
    1 & 2 & | & 0 \\
    3 & 4 & | & 1 \\
\end{bmatrix} \mapsto \begin{bmatrix}
    2 & 4 & | & 0 \\
    3 & 4 & | & 1 \\
\end{bmatrix} \]
This represents the row operation of scaling the first row by 2 ($2R_1 \mapsto R_1$). These matricies are equivalent but not equal.

\paragraph*{Example} Solve:
\begin{align*}
    &\begin{cases}
        2x + 3y - z = 3 \\
        -x - y + z = 1 \\
        x - 2y -5z = 0 \\
    \end{cases} = \begin{bmatrix}
        2 & 3 & -1 & | & 3 \\
        -1 & -1 & 1 & | & 1 \\
        1 & -2 & -5 & | & 0 \\
    \end{bmatrix} \xmapsto{R3 \mapsto R1} \begin{bmatrix}
        1 & -2 & -5 & | & 0 \\
        -1 & -1 & 1 & | & 1 \\
        2 & 3 & -1 & | & 3 \\
    \end{bmatrix} \xmapsto{R2 + R1 \mapsto R2} \begin{bmatrix}
        1 & -2 & -5 & | & 0 \\
        0 & -3 & -4 & | & 1 \\
        2 & 3 & -1 & | & 3 \\
    \end{bmatrix} \\
    &\xmapsto{R3 - 2R1 \mapsto R3} \begin{bmatrix}
        1 & -2 & -5 & | & 0 \\
        0 & -3 & -4 & | & 1 \\
        0 & 7 & 9 & | & 3 \\
    \end{bmatrix} \xmapsto{-\frac{1}{3}R2 \mapsto R2} \begin{bmatrix}
        1 & -2 & -5 & | & 0 \\
        0 & 1 & \frac{4}{3} & | & -\frac{1}{3} \\
        0 & 7 & 9 & | & 3 \\
    \end{bmatrix} \xmapsto{R3 - 7R2 \mapsto R3} \begin{bmatrix}
        1 & -2 & -5 & | & 0 \\
        0 & 1 & \frac{4}{3} & | & -\frac{1}{3} \\
        0 & 0 & \frac{-1}{3} & | & \frac{16}{3} \\
    \end{bmatrix} \\
    &\xmapsto{-3R3 \mapsto R3} \begin{bmatrix}
        1 & -2 & -5 & | & 0 \\
        0 & 1 & \frac{4}{3} & | & -\frac{1}{3} \\
        0 & 0 & 1 & | & -16 \\
    \end{bmatrix} \xmapsto{R2 - \frac{4}{3}R3 \mapsto R2} \begin{bmatrix}
        1 & -2 & -5 & | & 0 \\
        0 & 1 & 0 & | & 21 \\
        0 & 0 & 1 & | & -16 \\
    \end{bmatrix} \xmapsto{R1 + 5R3 \mapsto R1} \begin{bmatrix}
        1 & -2 & 0 & | & -80 \\
        0 & 1 & 0 & | & 21 \\
        0 & 0 & 1 & | & -16 \\
    \end{bmatrix} \\
    &\xmapsto{R1 + 2R2 \mapsto R1} \begin{bmatrix}
        1 & 0 & 0 & | & -38 \\
        0 & 1 & 0 & | & 21 \\
        0 & 0 & 1 & | & -16 \\
    \end{bmatrix} \implies (x, y, z) = (-38, 21, -16)
\end{align*}